\scanpage{153}
\mspage{724}

\unsure{ment}, \deleted{\illegible{1}} ou à expliciter en termes de ``fonctions elliptiques'' \unsure{standard}.

Sans m'attarder sur une telle explicitation, et revenant \inserted{plutôt} au plongement correspondant
\diagram{%
(\unsure{318})\qquad
\begin{tikzcd}[ampersand replacement=\&,column sep=large]
U \arrow[r,hook] \& \mathcal J^0_{U/S}
\end{tikzcd}}
et à la section \unsure{marquée}
\[
(\unsure{319})\qquad e\in U(S)=H^0(S,\mathcal J^0_{U/S}),
\]
\inserted{on peut utiliser la tour d'isogénies}
\deleted{on peut utiliser l'hom. canonique}
\diagram{%
\deleted{(\unsure{314})}\qquad
\begin{tikzcd}[ampersand replacement=\&,column sep=large]
\deleted{H^0(S,\mathcal J^0_{U/S})} \arrow[r,phantom,"\deleted{\longrightarrow}"{description}] \& \deleted{H}
\end{tikzcd}}
\diagram{%
(\unsure{321})\qquad
\begin{tikzcd}[ampersand replacement=\&,column sep=large]
1 \arrow[r,dash] \& {}_n\mathcal J^0_{U/S} \arrow[r,hook] \& \mathcal J^0_{U/S} \arrow[r,"n.\mathrm{id}_{\mathcal J^0_{U/S}}"] \& \mathcal J^0_{U/S} \arrow[r,dash] \& 1
\end{tikzcd}}
(pour $n$ variable), i.e. l'isogénie universelle déduite par passage à $\varprojlim$ :
\diagram{%
(\unsure{322})\qquad
\begin{tikzcd}[ampersand replacement=\&,column sep=large,row sep=large]
1 \arrow[r,dash] \& T_\infty(\mathcal J^0_{U/S}) \arrow[r,dash] \arrow[d,phantom,"\rotatebox{90}{$\simeq$}"{description}] \& \widetilde{\mathcal J}^0_{U/S} \arrow[r,"p"] \arrow[d,phantom,"\rotatebox{90}{$\underset{\mathrm{def}}{=}$}"{description}] \& \mathcal J^0_{U/S} \arrow[r] \& 1 \\
\& T_\infty(\mathbb G_m)_S\otimes_{\mathbb Z}\unsure{M_0}(S'/S) \& \varprojlim_n\unsure{{}_n}(\mathcal J^0_{U/S}) \& \&
\end{tikzcd}}
pour déduire, par image inverse de $e$, un torseur sous $T_\infty(\mathcal J^0_{U/S})\simeq T_\infty(\mathbb G_m)_S\otimes_{\mathbb Z}\unsure{M_0}(S'/S)$.

Ceci posé, on aura (cf (301) p. 719')
\[
(323)\qquad F_\infty(\widetilde X/S)\simeq p^{-1}(e),\quad\text{le }T_\infty(\mathcal J^0_{U/S})\text{-torseur}
\]
\inserted{déduit de la section $e$ de $\mathcal J^0_{U/S}$ sur $S$.}
\margin{152}

\scanpage{154}

\uline{Remarque} La classe d'isomorphie du $T_\infty(\mathcal J^0_{U/S})$-torseur obtenu est un élément de $H^1(S,T_\infty(\mathcal J^0_{U/S}))$, image de l'élément $e$ de $H^0(S,\mathcal J^0_{U/S})$, par le morphisme canonique
\diagram{%
(324)\qquad
\begin{tikzcd}[ampersand replacement=\&,column sep=large]
H^0(S,\mathcal J^0) \arrow[r] \& H^1(S,T_\infty(\mathcal J^0))
\end{tikzcd},}
dont le noyau est formé du sous-groupe de $H^0(S,\unsure{\mathcal J})$ des éléments ``infiniment divisibles''. Dans le cas où $S$ est \inserted{\unsure{réduit et est}} de type fini sur $\mathbb Z$ -- p.ex. \inserted{réduit} de type fini sur $\mathbb Q$ -- on sait (et c'est \deleted{\unsure{un théorème}} ``élémentaire'') que ce noyau est réduit à zéro. Cela signifie donc que la section $e$ \inserted{de $\unsure{U}$ sur $S$} (i.e. le ``paramètre elliptique'' envisagé) est alors entièrement déterminée, quand on connaît
\margin{153}

\scanpage{155}
\mspage{725}

$1^{\circ})$ le \deleted{\illegible{1}} \unsure{trirevêtement|trivêtement} étale $S'$ de $S$ associé, correspondant aux pts d'ordre $2$ de $\widetilde X/S$ -- d'où $\mathcal J^0\unsure{=}\mathcal J^0_{U/S}\simeq\mathcal J^0_0{}_{\deleted{\unsure{S}}}(\widetilde X,S')$ -- et

$2^{\circ})$ le $T_\infty(\mathcal J^{\unsure{0}}_{\deleted{\unsure{U/S}}})$-torseur $\unsure{\mathbb F_\infty|F_\infty}(\widetilde X/S)$ ;
\margin{154}

\scanpage{156}

$15^{\circ})$ Interprétation de $\Gamma_{\mathbb Q}\longrightarrow N_\sigma^{\unsure{\leftrightarrows}}$.

Au n° $11^{\circ}$ nous avons introduit l'hom. canonique
\diagram{%
\begin{tikzcd}[ampersand replacement=\&,column sep=huge,row sep=large]
\mathcal M\overset{\mathrm{def}}{=}\mathcal N_{1,1}^{\text{ét}} \arrow[r] \arrow[d,phantom,"\cup"{description}] \& \mathcal M_{\rho,\sigma}^{\sim} \arrow[d,phantom,"\cup"{description}] \\
\mathcal M^! \arrow[r] \& \mathcal M_{\rho,\sigma}^{\sim\unsure{\leftrightarrows}}
\end{tikzcd}\qquad\text{cf (217) p. \unsure{695})}}
que nous avons ``décomposé'' en \deleted{plus d'} deux homomorphismes (en plus d'un hom. trivial
\diagram{%
\begin{tikzcd}[ampersand replacement=\&,column sep=large,row sep=large]
\mathcal M \arrow[r] \arrow[dr] \& N_\rho^{\sim}\simeq\widehat{\mathbb Z}^{\times} \\
\& \Gamma_{\mathbb Q}^{\text{ét}} \arrow[u,"\chi"']
\end{tikzcd}}
) en \uline{deux} hom.
\diagram{%
(\unsure{325})\qquad
\begin{tikzcd}[ampersand replacement=\&,column sep=huge,row sep=large]
\mathcal M \arrow[r] \arrow[d,phantom,"\cup"{description}] \& G_{\rho,\sigma}\simeq\widehat{\mathbb Z}^{*}.\unsure{\mathcal G}^{+\wedge} \arrow[d,phantom,"\unsure{\cup}"{description}] \\
\mathcal M^! \arrow[r] \& G_{\rho,\sigma}^{\unsure{\leftrightarrows}}\simeq\widehat{\mathbb Z}^{*}.M^\wedge
\end{tikzcd}\qquad\text{(cf (218))}}
et
\diagram{%
(\unsure{326})\qquad
\begin{tikzcd}[ampersand replacement=\&,column sep=large]
\Gamma_{\mathbb Q}^{\text{ét}} \arrow[r] \& N_\sigma^{\unsure{\leftrightarrows}\deleted{\unsure{\sim}}}\simeq\widehat{\mathbb Z}^{*}.\widehat{\mathbb Z}
\end{tikzcd}\qquad\text{(cf (219)).}}

Le premier hom. (\unsure{322}) a été caractérisé en termes \inserted{\unsure{arithmético-}}\deleted{\unsure{algébriques}}géométriques (pour ce qui est de sa restriction à $\mathcal N_{1,1}\subset\mathcal N_{1,1}^{\text{ét}}$ lui-même, et en réduisant mod $\mu_G$ pour trouver un hom
\diagram{%
(327)\qquad\deleted{\unsure{$\mathcal M_{0,3}$}}\qquad
\begin{tikzcd}[ampersand replacement=\&,column sep=large]
\mathcal M_{0,3} \arrow[r,dash] \& G_{\rho,\sigma}^{\sim}\overset{\mathrm{def}}{\simeq}G_{\rho,\sigma}/\mu_G
\end{tikzcd},}
dans la section $13^{\circ})$, via la tour des
\margin{155}

\scanpage{157}
\mspage{\unsure{725|726}}

courbes de Fermat. Du point de vue actuel, où on essaye d'expliciter la géométrie algébrique de \inserted{\unsure{ess.\textsuperscript{t}}} type fini sur $\mathbb Q$ en termes d'\unsure{autom.} de groupes profinis, on peut dire aussi, inversement, que nous sommes arrivés à ``décrire'' la tour des courbes de Fermat, via l'hom (327) construit ici ``à la force du poignet''.

Il resterait maintenant à trouver une explicitation arithmético-géométrique de l'hom \deleted{\unsure{(322)}} de $\Gamma_{\mathbb Q}$ \inserted{$\subset\Gamma_{\mathbb Q}^{\text{ét}}$} dans $\widehat{\mathbb Z}^{*}.\widehat{\mathbb Z}$, induit par (\unsure{326}).
Mais \deleted{\unsure{cet l'hom.}} utilisant le plongement
(\diagram{%
\begin{tikzcd}[ampersand replacement=\&,column sep=large]
N_\sigma^{\unsure{\leftrightarrows}\deleted{\unsure{\sim}}} \arrow[r,hook] \& \unsure{G}_{\rho,\sigma}^{\unsure{\leftrightarrows}\deleted{\unsure{\sim}}}\simeq\widehat{\mathbb Z}^{*}.\uline{M}
\end{tikzcd},}
l'hom (\unsure{313}) apparaît comme le composé
\diagram{%
\begin{tikzcd}[ampersand replacement=\&,column sep=large]
\Gamma_{\mathbb Q}^{\text{ét}} \arrow[r,hook,"\unsure{\kappa}_{1/2}"'] \& \mathcal M^! \arrow[r,"(325)","\text{i.e. \textquotedblleft Fermat\textquotedblright}"'] \& G_{\rho,\sigma}^{\unsure{\leftrightarrows}\deleted{\unsure{\sim}}}
\end{tikzcd},}
de façon plus précise on a un diagr. comm. :
\margin{156}

\scanpage{158}
\[
\tag{328}
\diagram{\begin{tikzcd}[ampersand replacement=\&,column sep=large,row sep=large]
[\Gamma_{\mathbb Q}\hookrightarrow]\Gamma_{\mathbb Q}^{\text{ét}}
 \arrow[r] \arrow[d,"\kappa_{1/2}"'] \&
 \unsure{N|\mathcal N}_{\sigma}^{\unsure{\leftrightarrows|\sharp}}\simeq\widehat{\mathbb Z}^{*}.\widehat{\mathbb Z}
 \arrow[d,hook]\\
{}[\underbrace{\mathcal M_{0,3}^{!}}_{\Pi_1(\mathcal U_{0,3\mathbb Q})}\hookrightarrow]
 \mathcal M^{!\,\unsure{\infty|\text{ét}}}\simeq\unsure{N|\mathcal N}_{\unsure{\rho\beta|\rho3}}^{!\,\text{ét}}
 \arrow[r,"{\substack{\text{``Fermat''}\\(325)}}"'] \&
 G_{\rho,\sigma}^{\unsure{\leftrightarrows|\sharp}}\simeq\widehat{\mathbb Z}^{*}.\uline{M}^{\wedge}.
\end{tikzcd}}
\]

Cela signifie donc que l'on trouve l'hom. induit par (326)
\[
\tag{329}
\diagram{\begin{tikzcd}[ampersand replacement=\&]
\Gamma_{\mathbb Q}\arrow[r]\&\unsure{N|\mathcal N}_{\sigma}^{\unsure{\leftrightarrows|\sharp}}=\widehat{\mathbb Z}^{*}.\widehat{\mathbb Z}
\end{tikzcd}}
\]
en regardant l'image \inserted{dans $\unsure{N|\mathcal N}_{\sigma}^{\unsure{\leftrightarrows|\sharp}}\subset\widehat{\mathbb Z}^{*}.M^{\wedge}$} de l'hom.
\[
\tag{330}
\diagram{\begin{tikzcd}[ampersand replacement=\&]
\Gamma_{\mathbb Q}\arrow[r]\&\widehat{\mathbb Z}^{*}.\uline{M}^{\wedge}\subset\operatorname{Aff}(M^{\wedge}),
\end{tikzcd}}
\]
\deleted{dans $\unsure{N|\mathcal N}^{\unsure{\leftrightarrows|\sharp}}\subset\widehat{\mathbb Z}^{*}.\uline M$,}
déduit de la tour des courbes de fermat en regardant la fibre de la dite tour au point $1/2$, ce qui donne un torseur sous
$T_{\infty}(\mathbb G_m)\otimes_{\widehat{\mathbb Z}}\widehat M$,
s'explicitant galoisiennement justement par l'hom. (326).
Notons que le \deleted{\unsure{fibre}}\inserted{pt $1/2$}
\deleted{$U_{\infty}(1/2)$ est \unsure{isomorphe canoniquement},}
\inserted{de $\mathcal U_{0,3\mathbb Q}$ correspond, par l'iso. $\psi$ (317)}
de $\mathcal U_{0,3\mathbb Q}$ avec $U_{1,\mathbb Q}$, au point $(-\tfrac12,-\tfrac12,\unsure{1})$
\margin{157}

\scanpage{159}
\mspage{727}
de \deleted{\unsure{$U$}}
\[
\tag{331}
\inserted{s_{1/2}=(-\tfrac12,-\tfrac12,1)\in}
\deleted{\unsure{U_{1/2}}}
\left\{\begin{gathered}
U_{1\mathbb Q}\subset\mathbb G_{m\mathbb Q}^{3}/\mathbb G_{m\mathbb Q}\subset\mathbb P_{\mathbb Q}^{2}\\
\Vert\\
\{(x_0,x_1,x_{\infty})\in\Gamma(\mathbb G_{m\mathbb Q}^{3}/\mathbb G_{m\mathbb Q})\mid x_0+x_1+x_{\infty}=0\}
\end{gathered}\right.
\]
La fibre \inserted{\unsure{géo.}} de $U_n$ au ledit point s'identifie ainsi aux systèmes de solutions des équations ci-dessous:
\[
\tag{332}
\begin{aligned}
\deleted{\unsure{\bar s|s}}\quad U_n(s_{1/2})
&\simeq\{(x_0,x_1,x_{\infty})\in\Gamma\mathbb G_m^3\deleted{/\mathbb G_m}\mid x_0^n=-\tfrac12,\ x_1^n=-\tfrac12,\ x_{\infty}^{n}=\unsure{1|-1}\}\\[-2pt]
&\hspace{5em}\deleted{\unsure{\mathrm{mod}}\ \mathbb G_m}\\
&\simeq\{(\unsure{x_0},x_1,x_{\infty})\in\Gamma(\mathbb G_m^3)\mid x_0=-\tfrac12,\ x_1^n=-\tfrac12,\ x_{\infty}^n=\unsure{1|-1}\}\\
&\simeq\{(x_0,x_{\unsure{1|\infty}})\in\Gamma\mathbb G_m^2\mid x_0^n=\unsure{-\tfrac12},\ x_{\unsure{1|\infty}}^n=-\tfrac12\}\\
&\simeq\deleted{\unsure{\mu_n\times}}T_n(-\tfrac12)\times T_n(\unsure{\tfrac12|-\tfrac12})
\end{aligned}
\]
où on a posé, pour $\forall\ \unsure{\zeta|\xi}\in\mathbb G_m(S)$ ($S$ schém. quelc., \inserted{ici $S=\operatorname{Spec}\mathbb Q$})
\[
\tag{333}
\begin{gathered}
T_n(\unsure{\zeta|\xi})=
\deleted{\unsure{\text{image de}}\ \unsure{\zeta|\xi}\in H^0(S,\mathbb G_m)\ \text{par}}\\
\inserted{\mu_n\text{-torseur des racines }\deleted{\unsure{\text{de}}}n^{\text{ièmes}}\text{ de }\unsure{\zeta|\xi}},
\end{gathered}
\]

Passant à la limite on trouve un iso.
\[
\tag{334}
U_{\infty}(s_{1/2})\simeq T_{\infty}(\unsure{-\tfrac12})\times T_{\infty}(-\tfrac12),
\]
où on pose, comme de juste
\[
\tag{335}
T_{\infty}(\unsure{\zeta|\xi})=\varprojlim_n T_n(\unsure{\zeta|\xi})
\qquad(\text{un }T_{\infty}(\mathbb G_m)\text{-torseur}).
\]
\margin{158}

\scanpage{160}
Ainsi, l'hom (329) doit pouvoir s'interpréter en termes de l'hom galoisien
\[
\tag{\unsure{336|330}}
\diagram{\begin{tikzcd}[ampersand replacement=\&]
\Gamma_{\mathbb Q}\arrow[r]\&\operatorname{Gal}\bigl(\widetilde{\mathbb Q}[\sqrt[\infty]{-\tfrac12}]/\mathbb Q\bigr)
\end{tikzcd}}
\]
où $\widetilde{\mathbb Q}$ désigne l'extension cyclotomique universelle de $\mathbb Q$
\[
\tag{336}
\widetilde{\mathbb Q}\simeq\mathbb Q[\sqrt[\infty]{1}].
\]

Pour apprécier la ``taille'' de l'image de $\Gamma_{\mathbb Q}$ dans $\widehat{\mathbb Z}^{*}.\widehat{\mathbb Z}\simeq\unsure{N|\mathcal N}_{\sigma}^{\unsure{\leftrightarrows|\sharp}}$, il y a lieu de faire une digression de nature générale, sur les groupes de Galois \inserted{sur $\mathbb Q$} des extensions de la forme
\[
\widetilde{\mathbb Q}[\sqrt[\infty]{\unsure{\zeta_1|\xi_1}},\ldots,\sqrt[\infty]{\unsure{\zeta_N|\xi_N}}].
\]

Pour ceci, sur une base $S/\mathbb Q$, \unsure{on a pour simplifier} (on se restreindra au cas $S=\operatorname{Spec}\mathbb Q$) donnons nous un groupe abélien $A$ (\inserted{dans} le cas qui nous intéresse ici on prendra $A=\mathbb Z$), et
\margin{159}

\scanpage{161}
\mspage{728}
un hom
\[
\tag{337}
\diagram{\begin{tikzcd}[ampersand replacement=\&]
A\arrow[r,"\unsure{\zeta|\xi}"]\&\mathbb G_m(S)=H^0(S,\mathbb G_{mS})
\end{tikzcd}}
\]
(dans le cas présent, il est donné par $1\mapsto-1/2$),
ou ce qui revient au $\overline{\mathrm m}$, un hom.
\[
\tag{338\,\deleted{\unsure{bis}}}
\diagram{\begin{tikzcd}[ampersand replacement=\&]
A_S\arrow[r,"\unsure{\zeta|\xi}"]\&\mathbb G_{mS}.
\end{tikzcd}}
\]

Considérons la tour des isogénies de Kummer
\[
\tag{33\deleted{\unsure{8}}\inserted{\unsure{9}}}
\diagram{\begin{tikzcd}[ampersand replacement=\&]
1\arrow[r,dash]\&\mu_n\arrow[r,dash]\&\mathbb G_m\arrow[r,dash,"\lambda\mapsto\lambda^n"]\&\mathbb G_m\arrow[r,dash]\&1
\end{tikzcd}}
\]
d'où, par image inverse via (33\unsure{8}), une tour d'extensions de schémas en groupes étales sur $S$
\[
\tag{339}
\diagram{\begin{tikzcd}[ampersand replacement=\&]
1\arrow[r]\&\mu_n\arrow[r,dash]\&E_n(\unsure{\zeta|\xi})\arrow[r,dash]\&A_S\arrow[r,dash]\&1
\end{tikzcd}}
\]
et par \unsure{pass.} à la $\varprojlim$
\[
\tag{340}
\diagram{\begin{tikzcd}[ampersand replacement=\&]
1\arrow[r,dash]\&T_{\infty}(\mathbb G_{mS})\arrow[r,dash]\&E_{\infty}(\unsure{\zeta|\xi})\arrow[r,dash]\&A_S\arrow[r,dash]\&1.
\end{tikzcd}}
\]

Supposons $S$ \inserted{$0$-}connexe, et prenons un \deleted{\unsure{fixe}} pt géom \inserted{$\bar s$} de $S$, on aura donc un hom canonique
\[
\tag{341}
\diagram{\begin{tikzcd}[ampersand replacement=\&]
\Pi_1(S,\bar s)\arrow[r]\&\operatorname{Aut}_{\substack{\mathrm{gr}\\\text{pro-finis}}}\bigl(E_{\infty}(\unsure{\zeta|\xi})_{\bar s}\bigr),
\end{tikzcd}}
\]
\margin{160}

\scanpage{162}
où la fibre $E_{\infty}(\unsure{\zeta|\xi})_{\bar s}$ est elle-même munie d'une structure d'extension abélienne
\[
\tag{342}
\diagram{\begin{tikzcd}[ampersand replacement=\&]
1\arrow[r,dash]\&\underbrace{T_{\infty}(\bar k^{*})}_{T}\arrow[r]\&E_{\infty}(\unsure{\zeta|\xi})_{\bar s}\arrow[r,dash]\&A\arrow[r,dash]\&1.
\end{tikzcd}}
\]

L'hom (341) est compatible avec cette structure d'extension, quand on fait opérer $\Pi_1(S,\bar s)$ trivialement sur $A$, et sur $T_{\infty}(\bar k^{*})$ \deleted{\unsure{au} par la}\inserted{($\widehat{\mathbb Z}$-module libre de rg $1$)} par multiplication par le caractère cyclotomique
\[
\tag{343}
\diagram{\begin{tikzcd}[ampersand replacement=\&]
\Pi_1(S,\bar s)\arrow[r]\&\Pi_1(\mathbb Q,\bar s)\arrow[r,"\chi"]\&\widehat{\mathbb Z}^{*}.
\end{tikzcd}}
\]

Soit \deleted{$G$ le groupe}
\[
\tag{344}
\begin{gathered}
G_{\unsure{\zeta|\xi}}\subset\operatorname{Aut}(E_{\infty}(\unsure{\zeta|\xi})_{\bar s})\\
\Vert\\
\left\{u\in\operatorname{Aut}\ \middle|
\begin{gathered}
u\text{ respecte la structure d'ext en}\\
\text{induisant }\operatorname{id}_A\ \deleted{\text{de}}\text{ --- et }\inserted{\text{dans }T_{\infty}(\bar k^{*})}\text{ une}\\
\text{homothétie par un élément}\\
\mu\in\widehat{\mathbb Z}^{*}\text{ bien déterminé}
\end{gathered}\right\}
\end{gathered}
\]
\inserted{on a donc un hom canonique}\inserted{$u\mapsto\mu$}
\deleted{un groupe ne dépend que de $A$}
\[
\diagram{\begin{tikzcd}[ampersand replacement=\&]
G\arrow[r]\&\widehat{\mathbb Z}^{*},
\end{tikzcd}}
\]
s'insérant dans un diagramme
\[
\tag{345}
\diagram{\begin{tikzcd}[ampersand replacement=\&]
1\arrow[r]\&\operatorname{Hom}_{\mathbb Z}(A,T_{\infty}(\bar k^{*}))\arrow[r]\&G_{\unsure{\zeta|\xi}}\arrow[r]\&\widehat{\mathbb Z}^{*},
\end{tikzcd}}
\]
\margin{161}

\scanpage{163}
\mspage{729}

où l'image \deleted{est} de $G_{\unsure{\zeta|\xi}}\longrightarrow\widehat{\mathbb Z}^{*}$ est formée des $\mu\in\widehat{\mathbb Z}^{*}$ tels que \inserted{l'homothétie,} $\mu_T$ ne change pas la classe de l'ext. (342) de $A$ par $T$. L'hom (341) peut donc s'interpréter comme s'insérant dans un diagramme
\[
\tag{346}
\diagram{\begin{tikzcd}[ampersand replacement=\&,column sep=large,row sep=large]
1\arrow[r] \& \operatorname{Hom}_{\mathbb Z}(A,T_\infty(\bar k^{*}))\arrow[r] \& G_{\unsure{\zeta|\xi}}\arrow[r] \& \widehat{\mathbb Z}^{*}\\
\&\& \Pi_1\arrow[u,"\gamma_{\unsure{\zeta|\xi}}"]\arrow[ur,dashed,"{\text{``\unsure{caractère} cyclotomique''}}"',"{\text{comm.}}" description] \&
\end{tikzcd}}
\]

\margin{NB.}
Si $S$ est \unsure{ess.} de type fini sur $\mathbb Q$, alors l'image de $\Pi_1\longrightarrow\widehat{\mathbb Z}^{*}$ est un s-groupe ouvert de $\widehat{\mathbb Z}^{*}$, donc il en est a fortiori ainsi de l'image de $G_{\unsure{\zeta|\xi}}\longrightarrow\widehat{\mathbb Z}^{*}$ --- qui sera $\overline{\mathrm m}$ surjectif si le car. cyclotomique sur $\Pi_1$ est surjectif, i.e. si \deleted{\unsure{on}} les seules racines de $1$ sur $S$ sont $\pm1$ (p. ex. si $S=\operatorname{Spec}\mathbb Q\ldots$).

Le diagramme (346) a \unsure{une} une propriété de fonctorialité évidente pour $A$ variable (disons), i.e. pour un

\margin{162}

\scanpage{164}

hom de groupes abéliens
\[
\tag{347}
\diagram{\begin{tikzcd}[ampersand replacement=\&]
A'\arrow[r,"j"] \& A
\end{tikzcd}},
\]
\deleted{l'extension}
donnant un composé
\[
\tag{348}
\diagram{\begin{tikzcd}[ampersand replacement=\&,column sep=large]
A'\arrow[r,"j"]\arrow[rr,bend right=25,"\unsure{\zeta|\xi}'"'] \& A\arrow[r,"\unsure{\zeta|\xi}"] \& \mathbb G_m(S)
\end{tikzcd}},
\]
l'extension $E_\infty(\unsure{\zeta|\xi}')_{\bar s}$ s'identifie à l'image inverse de l'ext. $E_\infty(\unsure{\zeta|\xi})_{\bar s}$ par (347), d'où un hom
\margin{\deleted{\unsure{$G_{\unsure{\zeta|\xi}}$|$\mathbb Q_{\unsure{\zeta|\xi}}$}}}
\[
\diagram{\begin{tikzcd}[ampersand replacement=\&]
G_{\unsure{\zeta|\xi}}\arrow[r] \& G_{\unsure{\zeta|\xi}'}
\end{tikzcd}}
\]
\margin{injectif si $j$ est surjectif}
s'insérant dans un diagramme commutatif

\deleted{\unsure{(349)}}
\[
\tag{349}
\diagram{\begin{tikzcd}[ampersand replacement=\&,column sep=large,row sep=large]
1\arrow[r] \& \operatorname{Hom}_{\mathbb Z}(A',T)\arrow[r] \& G_{\unsure{\zeta|\xi}'}\arrow[r] \& \widehat{\mathbb Z}^{*}\\
1\arrow[r] \& \operatorname{Hom}_{\mathbb Z}(A,T)\arrow[u]\arrow[r] \& G_{\unsure{\zeta|\xi}}\arrow[u]\arrow[r] \& \widehat{\mathbb Z}^{*}\arrow[u,"\unsure{1}"]\\
\&\& \Pi_1\arrow[u,"\gamma_{\unsure{\zeta|\xi}}"]\arrow[uu,bend left=35,"\gamma_{\unsure{\zeta|\xi}'}"]\arrow[ur,"{\text{car. cyclot.}}"'] \&
\end{tikzcd}}
\]

où la flèche sur les $\operatorname{Hom}_{\mathbb Z}$ est transposée de (347). Ceci montre que pour étudier \inserted{la structure de} l'image de $\gamma_{\unsure{\zeta|\xi}}$ (346), on est ramené, par fonctorialité, \unsure{à en|à}

\margin{163}

\scanpage{165}
\mspage{730}

étudier le cas où $\unsure{j|\zeta}$ est \uline{injectif}
\[
\tag{350}
\unsure{j|\zeta}:\quad
\diagram{\begin{tikzcd}[ampersand replacement=\&]
A\arrow[r,hook] \& \mathbb G_m(S)
\end{tikzcd}}.
\]

\deleted{Je suppose disons}
On peut d'ailleurs, ``diviser'' encore un peu le diagramme (346), en introduisant
\[
\tag{351}
\Pi_1^{+}=\Pi_1^{+}(S,\bar s)\overset{\mathrm{def}}=
\operatorname{Ker}\bigl(\Pi_1(S,\bar s)\xrightarrow{\mathrm{cyclot.}}\widehat{\mathbb Z}^{*}\bigr),
\]
et l'hom de suites exactes
\[
\tag{352}
\diagram{\begin{tikzcd}[ampersand replacement=\&,column sep=large,row sep=large]
1\arrow[r] \& \Pi_1^{+}\arrow[r]\arrow[d,"\gamma_{\unsure{\zeta|\xi}}^{+}"] \& \Pi_1\arrow[r,"\chi_\Pi"]\arrow[d,"\gamma_{\unsure{\zeta|\xi}}"] \& \widehat{\mathbb Z}^{*}\arrow[d,"\unsure{1}"]\\
1\arrow[r] \& \operatorname{Hom}_{\mathbb Z}(A,T)\arrow[r] \& \deleted{\unsure{H|K}}\inserted{G}_{\unsure{\zeta|\xi}}\arrow[r] \& \widehat{\mathbb Z}^{*}.
\end{tikzcd}}
\]

On peut considérer que l'étude de la ``taille'' de $\operatorname{Im}(\gamma_{\unsure{\zeta|\xi}}:\Pi_1\longrightarrow G_{\unsure{\zeta|\xi}})$ est \unsure{essentiellement} ramenée à celle de la taille de $\operatorname{Im}\chi_\Pi$ et de celle de $\operatorname{Im}.\gamma_{\unsure{\zeta|\xi}}^{+}$; de façon précise, si on pose
\[
\tag{353}
\operatorname{Im}\gamma_{\unsure{\zeta|\xi}}=\widetilde\Pi_1,\quad
\operatorname{Im}\gamma_{\unsure{\zeta|\xi}}^{+}=\widetilde\Pi_1^{+},\quad
\operatorname{Im}\chi_\Pi=\deleted{\unsure{\widehat{\mathbb Z}^{*}|\widetilde\Gamma}}\inserted{\Gamma}
\]
on aura une suite exacte
\[
\tag{354}
\diagram{\begin{tikzcd}[ampersand replacement=\&]
1\arrow[r] \& \widetilde\Pi_1^{+}\arrow[r] \& \widetilde\Pi_1\arrow[r] \& \Gamma\arrow[r] \& 1
\end{tikzcd}}.
\]
Je voudrais essayer de donner des

\margin{164}

\scanpage{166}

indications sur la taille de $\widetilde\Pi_1^{+}\underset{\unsure{\zeta|\gamma}}{\subset}\underbrace{\operatorname{Hom}(A,T)}_{\substack{\simeq\operatorname{Hom}(A,\widehat{\mathbb Z})\\\deleted{\illegible{1}}\inserted{\unsure{\text{non can.}^{t}|\text{non can!}}}}}$ dans le cas où on suppose
\[
\tag{355}
\left\{\begin{array}{ll}
1^\circ) & A\quad\mathbb Z\text{-module de type fini}\\
2^\circ) & S\quad\text{ess.}^{\unsure{t|l}}\text{ de type fini sur }\mathbb Q\text{ et \uline{normal}}\\
3^\circ) & \unsure{j|\zeta}:A\longrightarrow\mathbb G_m(S)\text{ est injectif }\deleted{\unsure{\text{pur}}}.
\end{array}\right.
\]

L'hyp. $2^\circ)$ implique déjà on l'a vu, que $\Gamma\subset\widehat{\mathbb Z}^{*}$ est un sous-groupe ouvert. L'hyp. que $S$ est normal implique d'ailleurs que les groupes cherchés $\widetilde\Pi_1^{+}$ et $\Gamma$ ne changent pas \inserted{\unsure{en$^{h}$|en$^{t}$}} en remplaçant $S$ par son pt générique (on suppose pour simplifier que $\bar s$ est au dessus dudit pt générique), donc \unsure{sps} que $S=\operatorname{Spec}K$, où $K$ est un \inserted{corps} extension de type fini de $\mathbb Q$. Et on peut également supposer, si on y tient, que $S$ est affine, et lisse sur $\unsure{\mathbb Q}\ldots$.

Notons aussi, posant
\[
\tag{356}
\check A=\operatorname{Hom}_{\mathbb Z}(A,\mathbb Z),
\]

\margin{165}

\scanpage{167}
\mspage{731}

que $\check A$ est un $\mathbb Z$-module libre de t.f. (de rang \inserted{$n$} égal à $\operatorname{rg}_{\mathbb Z}A$), et \unsure{qu'on a} un isom canonique
\[
\tag{357}
\operatorname{Hom}_{\mathbb Z}\left(A,\underbrace{T}_{\substack{\simeq\widehat{\mathbb Z}\\(\unsure{\text{non can.}^{t}|\text{non can!}})}}\right)
\simeq
\underbrace{\check A\otimes_{\mathbb Z}T}_{\check A(-1)}
\quad
\left(\begin{gathered}
\simeq\widehat{\mathbb Z}^{n}\\
\unsure{\text{non canon.}^{t}|\text{non canon!}}
\end{gathered}\right)
\]

J'aurais envie de penser que l'image \deleted{\unsure{de}} $\widetilde\Pi_1^{+}$ de $\Pi_1$ dans $\operatorname{Hom}_{\mathbb Z}(A,T)\simeq\check A\otimes_{\mathbb Z}T$ est \uline{ouverte}. Je me demande $\overline{\mathrm m}$, dans le cas $S_1$ normal \inserted{conn.} de type fini sur $\mathbb Z$, \inserted{dimension $\unsure{2|\geq2}$} et $A=H^0(S_1,\mathbb G_m)$, pour $S=S_{1\mathbb Q}$, si l'hom
\[
\diagram{\begin{tikzcd}[ampersand replacement=\&,column sep=large]
\Pi_1=\Pi_1(S,\bar s)\arrow[r,no head] \& \operatorname{Hom}(A,T)\simeq\check A(-1)
\end{tikzcd}}
\]
n'est pas nécessairement surjectif??

(Plus \unsure{généralement}, il en est peut être ainsi si $\mathbb G_m(S)/\operatorname{Im}A$ est sans torsion, cas auquel on doit pouvoir tjrs se réduire, dans l'hyp. (355), quitte à agrandir $A$ en $A'\inserted{\subset H^0(S,\mathbb G_m)}$, \unsure{en|ou} un $A'$ tel que $A'/A$ soit fini\ldots).

\margin{166}

\scanpage{168}

Je vais revenir au cas $S=\operatorname{Spec}\mathbb Q$,
alors \deleted{\unsure{$\operatorname{Gal}$}}
\[
\mathbb G_m(\unsure{S})=\mathbb G_m(\mathbb Q)=\mathbb Q^*
\simeq\{\pm1\}\times\mathbb Z^{(\mathbb P)}
\]
où $\mathbb P$ est l'\unsure{ens} des \unsure{nbs} premiers,
\[
\unsure{\Pi_1|\pi_1}=\operatorname{Gal}(\overline{\mathbb Q}/\mathbb Q)=\Gamma_{\mathbb Q}.
\]

La question est donc d'étudier les hom
\[
\tag{358}
\unsure{\xi|\zeta}_p:\Gamma_{\mathbb Q}^{+}\longrightarrow\widehat{\mathbb Z}(-1)
\]
associés aux $p\in\mathbb P$, et de répondre notamment
aux questions

a) les $\unsure{\xi|\zeta}_p$ sont-ils surjectifs, ou du moins l'image
est-elle un sous-groupe ouvert ?

b) les $\unsure{\xi|\zeta}_p$ sont ils linéairement indépendants
sur $\widehat{\mathbb Z}$ ? Plus précisément, si $p_1\ldots p_n$ sont
des \unsure{nbs} premiers distincts, l'hom
\diagram{\begin{tikzcd}[ampersand replacement=\&,column sep=large]
\Gamma_{\mathbb Q}^{+}\arrow[r,"{(\unsure{\xi|\zeta}_{p_1},\ldots,\unsure{\xi|\zeta}_{p_n})}"']
\&\widehat{\mathbb Z}^{n}(\unsure{-1|1})
\end{tikzcd}}
est-il surjectif, ou du moins l'image de
$\Gamma_{\mathbb Q}^{+}$ est-elle un sous-groupe ouvert ?

Je vais commencer par regarder un peu
la question a) (le cas auquel nous en
\unsure{étions} dans le \unsure{contexte} de Fermat était
celui où $p=2$). Pour ceci, on décompose
$\unsure{\xi|\zeta}_p$ en ses composantes
\[
\tag{359}
\unsure{\xi|\zeta}_{p,\deleted{\unsure{\ell}}\ell}:\Gamma_{\mathbb Q}^{+}
\longrightarrow\widehat{\mathbb Z}_{\ell}(\unsure{-1|1})
\qquad(\ell\in\mathbb P)
\]

\margin{167}

\scanpage{169}
\mspage{732}

J'aimerais prouver donc \inserted{(au moins !)} que celles-ci sont
des hom. surjectifs, ou du moins \unsure{surjectifs}
sauf pour un \unsure{nb} fini de $\ell$, et à image
ouverte pour $\forall\ell$. Étudier $\unsure{\xi|\zeta}_{p,\ell}$, c'est
se restreindre à la tour des isogénies
\[
\tag{339}
\diagram{\begin{tikzcd}[ampersand replacement=\&,column sep=large]
1\arrow[r,dash]\&\mu_n\arrow[r,dash]\&\mathbb G_m
\arrow[r,"{\lambda\mapsto\lambda^n}"]\&\mathbb G_m\arrow[r,dash]\&1
\end{tikzcd}}
\]
pour $n$ puissance de $\ell$
\[
\tag{360}
n=\ell^\nu,
\]
et étudier \deleted{\unsure{l'|les}}\inserted{la tour des}
images inverses de $p\in \mathbb G_m(\mathbb Q)$
par ses isogénies. Il est naturel alors
de regarder ``le'' groupe de décomposition
\margin{\deleted{\unsure{(360)|(359)}}}
\[
\tag{361}
D_\ell\unsure{=|\simeq}\unsure{\Pi_1|\pi_1}(\mathbb Q_\ell,\overline{\mathbb Q}_\ell)
\hookrightarrow\Gamma_{\mathbb Q}
\]
et le caractère cyclotomique \inserted{$\ell$-adique} induit
\[
\tag{362}
\diagram{\begin{tikzcd}[ampersand replacement=\&,column sep=large]
D_\ell\arrow[r,"\chi_\ell"]\&\mathbb Z_\ell^*
\end{tikzcd}},
\]
ou mieux, le sous-groupe d'inertie
\[
\tag{363}
I_\ell\subset D_\ell
\]
et le caractère induit par $\chi_\ell$ sur $I_\ell$.
On aura en effet
\[
\tag{364}
I_\ell^+\overset{\mathrm{def}}{=}
\operatorname{Ker}(\chi_\ell:I_\ell\longrightarrow\mathbb Z_\ell^*)
\subset\Gamma_{\mathbb Q}^+
\]
car les $\chi_{\ell'}$ $(\ell'\ne\ell)$ sont triviaux sur $I_\ell$
(ils ne sont pas triviaux cependant sur $D_\ell$ !).

\margin{168}

\scanpage{170}

On va donc regarder la restriction
\[
\tag{365}
\unsure{\xi|\zeta}_{p,\ell}\mid I_\ell^+:
I_\ell^+\longrightarrow\widehat{\mathbb Z}_\ell(-1)
\]
et on peut se demander déjà si cet
homomorphisme est surjectif, ou si
du moins son image est \unsure{ouverte|ouvert} ?

\margin{$\unsure{=|-}$}

\margin{169}

\scanpage{171}
\mspage{733}

16°) \uline{Réflexions sur le théorème de Fermat.}

Le ``théorème de Fermat'' peut s'énoncer en
disant que pour $n\geq3$, la courbe de Fermat
affine sur $\mathbb Q$
\[
\tag{366}
\begin{aligned}
U_n\overset{\mathrm{def}}{=}
&\left\{(x_0,x_1,x_\infty)^{\cdot}\in\Gamma\mathbb P^2\ \middle|\
\begin{gathered}
x_0,x_1,x_\infty\ \mathrm{inv.}\\
x_0^n+x_1^n+x_\infty^n=0
\end{gathered}\right\}\\
\simeq{}&\left\{(x_0,\unsure{x_1|a_1},\unsure{x_\infty|a_\infty})^{\cdot}\in
\Gamma \mathbb G_m^{\deleted{\unsure{3}}\inserted{\{0,1,\infty\}}}/\mathbb G_m
\ \middle|\ x_0^n+x_1^n+x_\infty^n=0\right\}
\end{aligned}
\]
n'a pas de \deleted{\unsure{solution|solutions}} pt rat $/\mathbb Q$.
Nous avions introduit \deleted{\unsure{cette}}\inserted{la} ``tour'' des courbes
de Fermat dans §11°) (p. 696 ff), et dans
14°) est apparue la description de $U_n$ comme
image inverse de
\[
\tag{367}
\begin{array}{l}
U_1\subset\overbrace{\mathbb G_{m_{\mathbb Q}}\otimes_{\mathbb Z}M}^{T}
\simeq \mathbb G_{m_{\mathbb Q}}^{\{0,1,\infty\}}/\mathbb G_{m_{\mathbb Q}}\\[-.25em]
\quad\Vert\\[-.25em]
\left\{(x_0,x_1,x_\infty)^{\cdot}\in\Gamma \mathbb G_m^{\{0,1,\infty\}}/\mathbb G_m
\ \middle|\ x_0+x_1+x_\infty=0\right\}
\end{array}
\]
par l'isogénie de Kummer d'indice $n$
\[
\tag{368}
\diagram{\begin{tikzcd}[ampersand replacement=\&,column sep=large]
1\arrow[r]\&\mu_n\otimes_{\mathbb Z}M\arrow[r,dash]\&T
\arrow[r,"{x\mapsto x^n}"]\&T\arrow[r]\&1
\end{tikzcd}}.
\]

Calculons le genre \inserted{$g_n$} de $\widehat U_n$, qui est un
revêtement de degré $n^2$ de
$\widehat U_1\underset{\unsure{*|\ne}}{\simeq}
\widehat{\unsure{\mathcal U|\mathcal M}\deleted{\unsure{\mathbb Q|0,3}}}_{0,3}$,
\inserted{avec}\deleted{\unsure{totalement}}
ramification\deleted{dans}\inserted{de degré $n$ \unsure{en}}
les pts \inserted{au dessus de} $0,1,\infty$, et
étale en dehors. La formule de Hurwitz

\margin{170}
